\section{Introduction}

\subsection{Présentation générale}

S.Y.S.T.U.S. signifie \textit{Sound-Yielded System for Tracking \& User Sync}.
Le projet consiste à concevoir un objet musical intelligent, inspiré par
l'esthétique d'une cassette futuriste, capable d'écouter un extrait sonore,
d'identifier le morceau joué et d'afficher les informations principales sur un
écran e-paper.

Le prototype a été développé dans le cadre de l'option prototypage, sur une
période courte allant du 11 mai 2026 au 11 juin 2026. L'objectif n'est donc pas
de produire un objet industrialisé, mais de valider une chaîne fonctionnelle
complète : interaction utilisateur, capture audio, reconnaissance musicale,
affichage physique et configuration Wi-Fi locale.

Le projet a été réalisé par Antonin Mazaury et Victor Petit.

\subsection{Problématique}

Les services de reconnaissance musicale sont généralement associés à des
applications mobiles. S.Y.S.T.U.S. explore une autre approche : intégrer cette
fonction dans un objet autonome, manipulable, documenté et reproductible dans un
contexte de fablab.

La problématique principale est la suivante :

\begin{quote}
Comment concevoir en un mois un prototype physique capable de reconnaître une
musique, d'afficher un résultat lisible et de rester suffisamment documenté pour
être compris, testé et amélioré ?
\end{quote}

\subsection{Objectifs pédagogiques}

Le projet répond aux attendus de l'option prototypage :

\begin{itemize}
    \item formaliser un cahier des charges ;
    \item réaliser une analyse fonctionnelle ;
    \item utiliser des outils de management de projet ;
    \item intégrer de la modélisation 3D et de la programmation ;
    \item utiliser une carte programmable de type Raspberry Pi ;
    \item concevoir un prototype fonctionnel ;
    \item produire une documentation technique exploitable ;
    \item préparer une soutenance avec démonstration.
\end{itemize}

\subsection{Périmètre du document}

Ce document regroupe les éléments écrits attendus pour le rendu du 10 juin 2026
:

\begin{itemize}
    \item cahier des charges ;
    \item analyse fonctionnelle ;
    \item outils de management de projet : WBS, planning, Kanban, risques ;
    \item documentation technique ;
    \item critères de validation et plan de tests ;
    \item bilan des livrables et limites connues.
\end{itemize}

Le poster de présentation est volontairement exclu de ce dossier, car il sera
réalisé séparément.

\subsection{Modalités de rendu}

Le rendu de l'option prototypage est ouvert du lundi 11 mai 2026 à 00:00 au
mercredi 10 juin 2026 à 17:00. Les documents doivent respecter le nommage
demandé :

\begin{itemize}
    \item \texttt{SYSTUS\_CDC.PDF} pour le cahier des charges ;
    \item \texttt{SYSTUS\_Poster.PDF} pour le poster de présentation ;
    \item \texttt{SYSTUS\_DocTechnique.PDF} pour la documentation technique.
\end{itemize}

Les éléments évaluables sont :

\begin{itemize}
    \item rédaction d'un cahier des charges ;
    \item réalisation d'une analyse fonctionnelle ;
    \item mise en place d'outils de management de projet ;
    \item conception de la documentation technique ;
    \item création d'un poster de présentation ;
    \item archives : CAO, programme, schémas et fichiers utiles.
\end{itemize}

Le livrable documentaire est à rendre en PDF et sur la plateforme FabManager.
